\chapter{前瞻：生态下一步与决策准则}
\label{ch:forward}

\section{当前生态位评估（\versiontag{0.9.5} + Phase 12）}

\begin{table}[htbp]
\centering
\caption{本土化成功度（自评）}
\begin{tabular}{@{}llp{6cm}@{}}
\toprule
维度 & 状态 & 说明 \\
\midrule
本地内核可信 & 强 & 363+ gtest + commit-point + RAR/Merkle；ABI v29 封顶 \\
单文件吞吐 & 中上 & L5 $\sim$171；距 Stage 3.2 远 \\
多文件吞吐 & 强 & L7 $\sim$646+ \\
运维闭环 & 强 & snapshot/compact/repo-stats/daemon/service \\
云与多节点 & 外圈 & \filepath{sync\_cpp/} + S3；内核零网络 \\
流程成熟度 & 异类 & bench 驱动；非 enterprise 流程 \\
\bottomrule
\end{tabular}
\end{table}

\section{Phase 12 外侧云生态（已规划/落地）}

\begin{enumerate}
  \item \textbf{内核冻结}：ABI v29 封顶；\filepath{CLOUD\_ECOSYSTEM.md} 书面边界；
  \item \textbf{轨道 A}：EBB v2 delta 摆渡 + Workbench Sync Activity；
  \item \textbf{轨道 B}：\filepath{sync\_cpp/} + \texttt{eb-sync} + S3 SigV4；daemon \texttt{sync\_drain} 子进程；
  \item \textbf{不做}：ABI v30 云 API、FSM 内嵌 remote、双向多设备。
\end{enumerate}

\section{我将坚持的决策准则}

\begin{philobox}[Forward 准则]
\begin{enumerate}
  \item \textbf{Immutable floor 永不交易}——哪怕 L5 差 10 MB/s；
  \item \textbf{默认路径必须是 bench 最优}——实验路径 opt-in；
  \item \textbf{parity 失败 = 回退}，不「差不多」；
  \item \textbf{P5 超预算则合并拓扑}，不叠床架屋；
  \item \textbf{CHANGELOG 先于 slide}——生态对话用数字；
  \item \textbf{不恢复无用流程}——review 靠测试与文档审计。
\end{enumerate}
\end{philobox}

\section{给未来协作者}

若你加入维护：
\begin{itemize}
  \item 先读 \textbf{本文} + \filepath{PERF\_BASELINE.md}，再读 \filepath{kernel-manual}；
  \item 任何 pipeline/CDC 改动：先问 P1 还是 P3，再跑 parity；
  \item 若提议「恢复 Scrum/长设计评审」：请证明它比 \texttt{ebbackup\_bench\_check} 更便宜且更准。
\end{itemize}

\section{结语}

ebbackup 的本土化，不是把某套国外方法论翻译成中文，而是在\textbf{Windows 桌面、单人带宽、本地 MB/s 话语体系、可验证备份}的交叉压力下，长出一套\textbf{以测量为立法、以测试为执法、以剪枝为风控}的内核进化史。

工程手册记录「是什么」；本文记录「为什么是这样」。
